\documentclass[11pt,a4paper]{article}
\newcommand{\labonaamwaarde}{lab 3}
\usepackage[english]{syllabus}

\begin{document}

\syllabustitel{lab 3 -- \emph{Data Processing I}}{From file to tidy table}

% ---------------------------------------------------------------------------
% Shared by the syllabi of labs 3, 4 and 5. Edit it here, not in the lab files.
\section{Overview}

In labs 3, 4 and 5 you are the new data analyst of a (fictional) youth
football and athletics club. The head coach gives you three export files
from the club's systems. They are messy on purpose, the way real exports
are: inconsistent labels, sprint times in milliseconds, missing values,
an athlete who appears twice and a column with a comma as the decimal
separator.

Each session builds on the previous one:

\begin{description}[leftmargin=1.6cm, style=nextline]
  \item[Lab 3] \textbf{From file to tidy table.} Load, explore, filter,
    add columns and clean until you can trust the data.
  \item[Lab 4] \textbf{Combine, summarize, reshape.} Summarize groups,
    combine tables, switch between a wide and a long shape, and work with
    a heart-rate series per second.
  \item[Lab 5] \textbf{Tell the story with plotly.} Choose the right chart
    for a question, make honest charts with \code{plotly.express} and
    combine them into a small interactive dashboard.
\end{description}


% ---------------------------------------------------------------------------
\section{Learning outcomes}

After lab 3 you can:

\begin{enumerate}[label=\textbf{LO\arabic*}, leftmargin=1.4cm]
  \item load a CSV file with \code{pandas} and describe its structure,
        data types and missing values;
  \item select and filter rows and columns with boolean conditions, and
        sort the result;
  \item derive new columns (age, speed, BMI) and convert text to numbers
        and dates;
  \item clean a dataset and justify each choice: inconsistent text
        labels, wrong units, missing values and duplicates.
\end{enumerate}

% ---------------------------------------------------------------------------
\section{Prerequisites}

Labs 1 and 2 (\emph{Introduction to Python I and II}): variables and data
types, operators, \code{if}/\code{elif}/\code{else}, functions, string
methods, lists, dictionaries, loops, and reading a CSV file with the
\code{csv} module. Lab 2 ends by summarizing \code{athletes.csv} with
plain Python; in lab 3 you see how \code{pandas} does the same in a few
lines.

% ---------------------------------------------------------------------------
% Shared by the syllabi of labs 3, 4 and 5. Edit it here, not in the lab files.
\section{Software and materials}

Every lab is a \textbf{marimo notebook} in the \code{labs/} folder of the
course repository. You can open it in two ways:

\begin{itemize}
    \item \textbf{Online:} go to \url{https://marimo.app}, choose
        \emph{New $\rightarrow$ Open from URL\dots} and paste the raw URL
        of the notebook on GitHub. The data is then read from GitHub, so
        you don't have to install anything.
    \item \textbf{Locally:} install the packages from
        \code{requirements.txt} and open the notebook with marimo. The
        data is then read from the \code{data/} folder.
\end{itemize}

Three things to know about marimo:

\begin{itemize}
  \item Notebooks are \textbf{reactive}: when you change a cell, every
        cell that depends on it runs again, in the order of their
        dependencies.
  \item \textbf{Two cells cannot define the same variable.} Give every
        step its own clear name, such as \code{tests\_raw} and
        \code{tests\_clean}, or keep a sequence of steps together in one
        cell.
  \item Every exercise contains a \code{...} placeholder. A hidden cell
        below it checks your answer: it turns \textbf{green} when the
        answer is right, and otherwise gives a general hint about what is
        still wrong, without giving the answer away.
\end{itemize}


% ---------------------------------------------------------------------------
% Shared by the syllabi of labs 3, 4 and 5. Edit it here, not in the lab files.
\section{The dataset}

\noindent
\rowcolors{2}{tint}{white}
\begin{tabularx}{\textwidth}{@{}L{3.1cm}L{3.2cm}Y@{}}
\toprule
\textbf{File} & \textbf{Size} & \textbf{Content and known problems} \\
\midrule
\code{athletes.csv}   & 43 rows, 7 columns   & ID, name, birthdate, gender, position, height, weight. 13 spellings for 4 positions, weight as text with a decimal comma (\code{"57,5"}), one athlete entered twice with two IDs, a few missing values. \\
\code{tests.csv}      & 114 rows, 8 columns  & 10\,m and 30\,m sprint, countermovement jump height and Yo-Yo level, at 3 test moments. Some sprint times are stored in milliseconds; a few missing values. \\
\code{hr\_session.csv} & 12,000 rows, 5 columns & Heart rate and speed per second for 5 athletes during one 40-minute training session. \\
\bottomrule
\end{tabularx}
\rowcolors{1}{}{}


% ---------------------------------------------------------------------------
\needspace{18\baselineskip}
\section{Schedule}

% Schedule of one lab. Used by this lab's syllabus and by the syllabus for labs 3-5.
\textbf{Goal:} load and explore a dataset, select and filter rows and
columns, add new columns and clean the data.

\begin{schedule}
13:00--13:20 & Warm-up: why not just Excel? & Look at the raw text of a CSV file and see what a table library adds. \\
13:20--14:00 & Loading and exploring & \code{read\_csv}, \code{.head()}, \code{.shape}, \code{.dtypes}, \code{.info()}, \code{.describe()}, \code{mo.ui.table}. \\
14:00--14:50 & Selecting and filtering & Selecting columns, \code{.loc} with boolean conditions, combining conditions, \code{sort\_values}. \\
14:50--15:05 & Break & \\
15:05--15:50 & New columns & Text to number, age from a birthdate, speed from a sprint time, BMI. \\
15:50--17:00 & Cleaning & Missing values (\code{dropna} or \code{fillna}), text labels, units, duplicates, one complete cleaning pipeline. \\
17:00--17:30 & Challenge & The 5 fastest athletes per position. Exit ticket. \\
\end{schedule}

\textbf{Exercises:} (1) load \code{tests.csv}; (2) forwards younger than
16; (3) suspicious sprint times; (4) BMI; (5) age in years; (6) clean the
positions; (7) fix the sprint units; (8) remove the duplicate athlete;
\textbf{challenge:} the 5 fastest per position.


% ---------------------------------------------------------------------------
\section{Working method}

\begin{itemize}
  \item Each block starts with a short demo in the notebook, followed by
        exercises you solve in the notebook itself.
  \item Use the automatic checks as feedback, not as the goal: when a
        check turns green, make sure you can explain \emph{why} your code
        works.
  \item The lab ends with an exit ticket: one technique you will reuse
        and one thing that is still unclear. Bring the unclear points to
        lab 4.
\end{itemize}

% ---------------------------------------------------------------------------
\section{Preparation}

\begin{itemize}
  \item Check before the session that you can open the notebook, locally
        or on marimo.app.
  \item Review lists, dictionaries and functions from lab 2: pandas uses
        them all the time.
  \item Read the course text for lab 3 next to the notebook.
\end{itemize}

\end{document}
